\section*{Chapitre \ref{ch:b1:gene-expression} --- L'expression génétique~: transcription et traduction}

\begin{solution}{exo:b1:gene-expression:1}
\omterm{def:b1:nucleic-acids:chain}{ARN} $5'$-AUGCCUAAG-$3'$~; brin codant $5'$-ATGCCTAAG-$3'$.
\end{solution}

\begin{solution}{exo:b1:gene-expression:2}
Une coiffe en $5'$ (guanine modifiée), une \omterm{prop:b1:gene-expression:processing}{queue poly-A} en $3'$, et le
retrait des \omterm{def:b1:genomes:gene}{introns} par épissage.
\end{solution}

\begin{solution}{exo:b1:gene-expression:3}
AUG CCG AAA GUU UGA~: Met-Pro-Lys-Val, puis stop.
\end{solution}

\begin{solution}{exo:b1:gene-expression:4}
A~: l'ARNt chargé entrant est admis et vérifié. P~: l'ARNt qui porte la
chaîne en croissance~; la liaison se forme entre cette chaîne et l'\omterm{def:b1:proteins:aminoacid}{acide
aminé} du site A. E~: l'ARNt vidé, en partance.
\end{solution}

\begin{solution}{exo:b1:gene-expression:5}
Codant $2400 - 600 = 1800$ \omterm{def:b1:nucleic-acids:nucleotide}{nucléotides}~: 599 résidus (600 \omterm{def:b1:gene-expression:code}{codons}, dont
un stop). Durée $599/4 = \qty{150}{s}$. \omterm{def:b1:gene-expression:ribosome}{Ribosomes} $1800/90 = 20$ à la
fois.
\end{solution}

\begin{solution}{exo:b1:gene-expression:6}
Le message est lu par trois à partir d'un départ fixé~; une ou deux
bases en trop décalent tous les \omterm{def:b1:gene-expression:code}{codons} suivants et le reste de la
\omterm{def:b1:proteins:peptide}{protéine} n'est plus que du charabia, qui se termine d'ordinaire à un
stop prématuré. Trois bases en trop ajoutent un \omterm{def:b1:gene-expression:code}{codon} et rétablissent le
cadre~: la \omterm{def:b1:proteins:peptide}{protéine} a un résidu de plus et quelques résidus faux entre
les insertions, et elle fonctionne souvent encore.
\end{solution}

\begin{solution}{exo:b1:gene-expression:7}
CAG (Gln) en UAG (stop)~: la chaîne s'arrête au résidu 150, \omterm{def:b1:proteins:peptide}{protéine}
tronquée et presque certainement inactive (mutation non-sens). CAG en
CAA~: toujours Gln, silencieuse. CAG en CGG~: Arg à la place de Gln, une
substitution (faux-sens) dont l'effet dépend de la position.
\end{solution}

\begin{solution}{exo:b1:gene-expression:8}
Le \omterm{def:b1:gene-expression:ribosome}{ribosome} ne vérifie que l'appariement codon--anticodon~; il ne voit
pas l'\omterm{def:b1:proteins:aminoacid}{acide aminé}. Si une cystéine fixée sur son ARNt est convertie
chimiquement en alanine, le \omterm{def:b1:gene-expression:ribosome}{ribosome} insère de l'alanine partout où le
message dit cystéine~: c'est l'ARNt, non l'\omterm{def:b1:proteins:aminoacid}{acide aminé}, qui est lu. La
synthétase, qui apparie chaque \omterm{def:b1:proteins:aminoacid}{acide aminé} au bon ARNt, est donc le vrai
traducteur.
\end{solution}

\begin{solution}{exo:b1:gene-expression:9}
$400\times 4 = 1600$ liaisons~; les deux GTP du \omterm{def:b1:gene-expression:ribosome}{ribosome} par résidu en
font la moitié, les deux des synthétases l'autre moitié.
\end{solution}

\begin{solution}{exo:b1:gene-expression:10}
Chez les eucaryotes, le transcrit est fait dans le noyau et les
\omterm{def:b1:gene-expression:ribosome}{ribosomes} sont dans le \omterm{def:b1:eukaryotic-cell:organelle}{cytosol}, séparés par l'enveloppe~; et le
transcrit n'est pas un messager tant qu'il n'a pas été épissé et coiffé.
Cette séparation rend possible la maturation elle-même --- l'épissage
alternatif, le contrôle de l'exportation, et la vérification que le
message est complet avant d'être lu.
\end{solution}

\begin{solution}{exo:b1:gene-expression:11}
Quatre messagers~: avec les deux \omterm{def:b1:genomes:gene}{exons} (190 \omterm{def:b1:nucleic-acids:nucleotide}{nucléotides} ajoutés), avec
le 3 seul (90), avec le 4 seul (100), avec aucun des deux (0). Le cadre
est conservé si la longueur ajoutée est un multiple de 3~: 0 et 90 le
conservent~; 100 et 190 le décalent, de sorte que les messagers à \omterm{def:b1:genomes:gene}{exon} 4
seul ou à deux \omterm{def:b1:genomes:gene}{exons} lisent l'\omterm{def:b1:genomes:gene}{exon} 5 hors cadre et tronquent la
\omterm{def:b1:proteins:peptide}{protéine}. L'épissage alternatif doit respecter le cadre, et c'est pour
cela que les \omterm{def:b1:genomes:gene}{exons} sont souvent des multiples de trois.
\end{solution}

\begin{solution}{exo:b1:gene-expression:12}
Arbitraire~: rien dans la chimie ne dit que UUU doive signifier Phe~;
les attributions sont des conventions tenues par les synthétases.
Optimisé~: la troisième position est la plus dégénérée, de sorte que les
erreurs et les mutations qui y tombent sont souvent silencieuses, et des
\omterm{def:b1:gene-expression:code}{codons} qui diffèrent d'une base codent souvent des \omterm{def:b1:proteins:aminoacid}{acides aminés}
voisins, de sorte qu'une substitution est souvent bénigne --- le code
minimise les dégâts de l'erreur. Figé~: une synthétase qui changerait de
spécificité altérerait d'un coup toutes les \omterm{def:b1:proteins:peptide}{protéines} de la \omterm{def:b1:cell-unit-of-life:cell}{cellule}, des
milliers à la fois~; presque aucune \omterm{def:b1:cell-unit-of-life:cell}{cellule} n'y survit, de sorte que le
code ne peut pas dériver et est resté fixé depuis l'ancêtre commun de
tous les êtres vivants.
\end{solution}

\begin{solution}{pb:b1:gene-expression:1}
\textbf{1.} $30\,000/30 = \qty{1000}{s}$, soit environ \qty{17}{min}.
\textbf{2.} $28\,000/30\,000 = \qty{93}{\%}$.
\textbf{3.} $30\,000/2000 = 15$.
\textbf{4.} $2\times 30\,000 = \num{60000}$ \omterm{def:b1:water-small-molecules:atp}{ATP}.
\textbf{5.} $1000/10 = 100$ polymérases sur le \omterm{def:b1:genomes:gene}{gène}~; 360 messagers par
heure.
\textbf{6.} Sept \omterm{def:b1:genomes:gene}{introns}, de longueur moyenne $28\,000/7 = \qty{4000}{}$
\omterm{def:b1:nucleic-acids:nucleotide}{nucléotides}.
\textbf{7.} La transcription (\qty{17}{min})~; les \omterm{def:b1:genomes:gene}{introns} sont épissés
à mesure qu'ils sont faits, et le dernier n'ajoute qu'une minute.
\textbf{8.} \omterm{def:b1:genomes:gene}{Gène} $30\,000\times 0.34\,\mathrm{nm} = \qty{10}{\micro m}$~; messager \qty{0.68}{\micro m}.
\textbf{9.} $500/4 = \qty{125}{s}$.
\textbf{10.} $2000/80 = 25$.
\textbf{11.} Une chaîne s'achève toutes les $80/4 = \qty{20}{s}$~: 180
par heure.
\textbf{12.} $2/\ln 2 = \qty{2.9}{h}$ de production pleine~: environ 520
\omterm{def:b1:proteins:peptide}{protéines}.
\textbf{13.} $4\times 500 = 2000$ \omterm{def:b1:water-small-molecules:atp}{ATP}.
\textbf{14.} $60\,000/520 = 115$ \omterm{def:b1:water-small-molecules:atp}{ATP} par \omterm{def:b1:proteins:peptide}{protéine}.
\textbf{15.} Environ 2100 \omterm{def:b1:water-small-molecules:atp}{ATP}, dont \qty{95}{\%} pour la traduction.
\textbf{16.} $1 - (1 - 10^{-5})^{1503} \approx 1503\times 10^{-5} =
\qty{1.5}{\%}$.
\textbf{17.} $1 - (1 - 10^{-4})^{500} \approx \qty{5}{\%}$.
\textbf{18.} Transcription~: $1.5\%\times 0.75\times 0.67 \approx
0.75\%$ de messagers défectueux, donc de \omterm{def:b1:proteins:peptide}{protéines}~; traduction~:
$5\%\times 0.67 = 3.3\%$~; environ \qty{4}{\%} des molécules.
\textbf{19.} Une erreur de \omterm{def:b1:proteins:peptide}{protéine} se limite à une molécule, bientôt
dégradée~; une erreur de messager est recopiée dans des centaines de
\omterm{def:b1:proteins:peptide}{protéines}~; une erreur de réplication est héritée par tous les
descendants à jamais. Le coût d'une erreur, et donc l'exactitude qui
vaut d'être payée, croît à chaque étape en remontant.
\textbf{20.} $3000/180 = 17$ messagers présents.
\textbf{21.} $17\times 25 = 420$ \omterm{def:b1:gene-expression:ribosome}{ribosomes}.
\textbf{22.} $\num{3e9}\times 4 = \num{1.2e10}$ \omterm{def:b1:water-small-molecules:atp}{ATP} par jour
$= \qty{2e-14}{mol}$, \qty{1e-9}{J} par jour, soit \qty{1.2e-14}{W}
--- \qty{6}{fW} par micromètre cube.
\textbf{23.} $10^9$ \omterm{def:b1:water-small-molecules:atp}{ATP} par seconde font $\num{8.6e13}$ par jour~: la
synthèse protéique en représente le centième dans cette \omterm{def:b1:cell-unit-of-life:cell}{cellule} à
renouvellement lent (chez une bactérie en division, la moitié).
\textbf{24.} Plus de messager mûr~: les transcrits primaires
s'accumulent dans le noyau et la \omterm{def:b1:proteins:peptide}{protéine} disparaît à mesure que ses
messagers se dégradent, en quelques heures. La \omterm{def:b1:proteins:peptide}{protéine} bactérienne,
dont le \omterm{def:b1:genomes:gene}{gène} n'a pas d'\omterm{def:b1:genomes:gene}{introns}, n'est pas touchée.
\textbf{25.} Environ 2100 \omterm{def:b1:water-small-molecules:atp}{ATP} par molécule --- 2000 pour la traduction,
une centaine pour sa part du transcrit~; première \omterm{def:b1:proteins:peptide}{protéine} après
\qty{1000}{s} de transcription, une minute de maturation et
d'exportation, et \qty{125}{s} de traduction~: environ vingt minutes.
\end{solution}
